\section{本章小结}
\label{sec:graph-electrical-networks-summary}

给定各顶点的外部注入后，逐边的电流守恒和Ohm定律联立为拉普拉斯方程$\bm L\bm v = \bm b$。正权连通分量内注入总量为零是有解的充要条件。接地选定电势基准——删去的那一列对应已固定的电势，删去的那一行可由其余方程恢复。非连通时各分量独立选基准，孤点注入必须为零。

同一组注入下，电流不唯一决定于守恒条件。不同的守恒流与实际电流之差是环流——它不改变任何顶点的净流出。叠加环流只增加该环流自身的耗散，交叉项因电势与环流正交而消去。实际电流因此是耗散最小的守恒分流，最小耗散同时等于两端压降。这把能量最小原理和有效电阻联系在一起。

计算方式取决于任务。查询单个点对时，解它的单位注入方程再取电势差。树上每条边的电流可以由一侧子树的注入和直接确定，电势由根向叶递推，时间为$O(n)$。若要查询大量点对，可以用一批随机右端给每个顶点预计算低维坐标，之后每次查询只需取两列的平方距离。投影维数和方程求解精度分别控制近似误差和失败概率。

有效电阻在稀疏化与行走中还会出现：第~\ref{chap:graph-sparsification}章用边的有效电阻决定抽样权重，第~\ref{chap:graph-walks}章用点对有效电阻计算随机游走的平均往返时间。
